% Generated from locale/tr/content/applied-modal-logic/epistemic-logic/bisimulations.tex; do not hand-edit.


\olfileid[tr]{aml}{el}{bsd}

\olsection{Bisimülasyonlar}

Epistemik dilimizin ifade gücü konusunda geriye kalan sorulardan biri,
modeller ile bu modellerde geçerli formüller arasındaki ilişkiyle ilgilidir.
Çerçeve karşılıklılığı sonuçlarımızdan, belirli formüller bir çerçevede
geçerliyse o çerçevenin belirli özellikleri sağlamasını da güvenceye
alacaklarını gördük. Fakat modal dilimiz, örneğin üzerinde yansımalı bir ok
bulunan bir dünya ile her biri bir sonrakine giden sonsuz bir dünyalar
zincirini ayırt etmemize izin verir mi? Başka deyişle, bu iki dünyadan
yalnızca birinde geçerli olabilecek herhangi bir $A$ formülü var mıdır?

Bisimülasyon, bağıntısal modeller arasında tanımlayarak modellerin fiilen aynı
yapıya sahip olduğunu söylememizi sağlayan bir bağıntıdır. Göreceğimiz gibi,
epistemik dilimizde yakalanabilen bir model denkliği anlayışını belirler.

\begin{defn}[Bisimülasyon]
$M_1 = \tuple{W_1, R_1, V_1}$ ve $M_2 = \tuple{W_2, R_2, V_2}$ iki
bağıntısal model, $\mathcal{R} \subseteq W_1 \times W_2$ de ikili bir bağıntı
olsun. $\mathcal{R}$ bağıntısına \emph{bisimülasyon} deriz; bunun için her
$\tuple{w_1, w_2} \in \mathcal{R}$ bakımından aşağıdakiler geçerli olmalıdır:

\begin{enumerate}
  \item $w_1 \in V_1(p)$ ancak ve ancak $w_2 \in V_2(p)$ ise; bu, bütün
  !!{propositional variable}s~$p$ için geçerlidir.
  \item Bütün $a \in G$ özneleri ve $v_1 \in W_1$ dünyaları için
  $R_{1_a} w_1 v_1$ ise, bazı $v_2 \in W_2$ bulunup $R_{2_a} w_2 v_2$ ve
  $\tuple{v_1, v_2} \in \mathcal{R}$ ise.
  \item Bütün $a \in G$ özneleri ve $v_2 \in W_2$ dünyaları için
   $R_{2_a} w_2 v_2$ ise, bazı $v_1 \in W_1$ bulunup $R_{1_a} w_1 v_1$ ve
   $\tuple{v_1, v_2} \in \mathcal{R}$ ise.
\end{enumerate}
% [tr source emendation] Upstream uses undeclared agent set `A` in the two
% forth/back clauses; the chapter declares and consistently uses `G`.

$M_1$ ile $M_2$ arasında $w_1$ ve $w_2$ dünyalarını bağlayan bir bisimülasyon
bulunduğunda $\tuple{M_1, w_1} \leftrightarroweq \tuple{M_2, w_2}$ de
yazabilir ve $\tuple{M_1, w_1}$ ile $\tuple{M_2, w_2}$'ye
\emph{bisimüler} diyebiliriz.
\end{defn}

Bisimülasyon bağıntısındaki farklı koşullar farklı şeyleri güvenceye alır. 1.
koşul, bütün !!{propositional variable}s üzerinde uyuşmayı güvenceye aldığı
için bisimüler dünyaların aynı modal içermeyen formülleri sağlamasını sağlar.
Sırasıyla ``ileri'' ve ``geri'' koşulları diye anılan diğer iki koşul ise
erişilebilirlik bağıntılarının aynı yapıya sahip olmasını güvenceye alır.

\begin{thm}
$\tuple{M_1, w_1} \leftrightarroweq \tuple{M_2, w_2}$ ise her
!!{formula}~$!A$ için $\mSat{M_1}{!A}[w_1]$ ancak ve ancak
$\mSat{M_2}{!A}[w_2]$ ise geçerlidir.
\end{thm}

\begin{figure}
  \begin{center}
    \begin{tikzpicture}[modal]
      \node[world] (w1) {$w_1$}; 
      \node[world] (w2) [above left=of w1]{$w_2$}; 
      \node[world] (w3) [above right=of w1] {$w_3$};
      \draw[<->] (w1) to node {$a$} (w2);
      \draw[->,loop below] (w1) to node {$a$} (w1);
      \draw[<->] (w1) to [swap] node {$a$} (w3);
      \draw[->,loop above] (w2) to node {$a$} (w2) ;
      \draw[->,loop above] (w3) to node {$a$} (w3) ;
      
      \node[world](v1)[right of=w1, xshift=1.5in]{$v_1$};
      \node[world](v2)[above of=v1]{$v_2$};
      \draw[<->] (v1) to node {$a$} (v2);
      \draw[->,loop below] (v1) to node {$a$} (v1);
      \draw[->,loop above] (v2) to node {$a$} (v2) ;

      \draw[-,dotted] (w1) to (v1) ;
      \draw[-,dotted,bend left=45] (w2) to (v2) ;
      \draw[-,dotted,bend left=30] (w3) to (v2) ;
    \end{tikzpicture}
  \end{center}
  \caption{İki bisimüler model.}
  \ollabel{fig:bisimilar}
\end{figure}

\olref{fig:bisimilar}'de gösterilen iki model tam olarak aynı olmasa da
$w_1$ ve~$v_1$ dünyalarını bağlayan bir bisimülasyon vardır. Bu bisimülasyon
hem $w_2$'yi hem $w_3$'ü de~$v_2$'ye bağlar; bunun ardındaki düşünce, modal
dilimizde ifade edilebilen hiçbir şeyin onları gerçekten ayıramamasıdır.
Ancak $w_2$ ile~$w_3$ farklı !!{propositional variable}s{}ni sağlasaydı durum
farklı olurdu.

